% Generated from src/content/cv*.json by scripts/render_cv.py.
\documentclass[11pt,a4paper]{article}
\usepackage{fontspec}
\setmainfont{TeX Gyre Heros}
\setsansfont{TeX Gyre Heros}
\setmonofont{Latin Modern Mono Light}[Scale=MatchLowercase]
\renewcommand{\familydefault}{\sfdefault}
\usepackage[top=0.6in,bottom=0.6in,left=0.7in,right=0.7in,headheight=14pt,headsep=12pt]{geometry}
\usepackage{xcolor,hyperref,needspace,fancyhdr,lastpage,fontawesome}
\definecolor{ink}{HTML}{171717}
\definecolor{muted}{HTML}{555555}
\definecolor{paper}{HTML}{F3F3F1}
\hypersetup{hidelinks,pdftitle={Ruşen Birben CV},pdfauthor={Ruşen Birben}}
\setlength{\parindent}{0pt}
\setlength{\parskip}{0pt}
\raggedright
\linespread{1.05}
\color{ink}
\pagestyle{fancy}
\fancyhf{}
\fancyhead[L]{\ttfamily{}\small{}Ruşen Birben\enspace{} / CV}
\fancyfoot[R]{\ttfamily{}\footnotesize{}\thepage\ / \pageref*{LastPage}}
\renewcommand{\headrulewidth}{0pt}
\fancypagestyle{firstpage}{\fancyhead{}}
\newcounter{cvsec}
\newcommand{\cvsection}[2]{%
  \par\needspace{7\baselineskip}\vspace{6pt}\refstepcounter{cvsec}%
  {\setlength{\fboxsep}{3pt}\colorbox{ink}{\color{white}\ttfamily{}\bfseries{}\small{}\ifnum\value{cvsec}<10 0\fi\thecvsec}}%
  \hspace{7pt}{\large{} #1}\hspace{7pt}{\ttfamily{}\bfseries{}\fontsize{12}{14}\selectfont{} #2}%
  \par\vspace{4pt}{\color{ink}\hrule height 0.9pt}\nopagebreak\vspace{7pt}%
}
\newcommand{\dates}[1]{\hfill{\ttfamily{}\small{}\color{muted}#1}}
\newcommand{\cvbody}[1]{{\small{}\raggedright #1\par}}
\newcommand{\extarrow}{\,{\scriptsize\faExternalLink}}
\begin{document}
\thispagestyle{firstpage}
\hrule height 2.5pt\vspace{6pt}
{\ttfamily{}\fontsize{30}{33}\selectfont{}\bfseries{} Ruşen Birben}\par\vspace{3pt}
{\ttfamily{}\fontsize{11}{13}\selectfont{}Yapay Zeka Araştırmacısı ve Mühendisi}\par\vspace{6pt}
{\small{}
\faMapMarker\enspace{}Ankara, Türkiye\quad{}
\faEnvelope\enspace{}\href{mailto:contact@rusen.ai}{contact@rusen.ai}\quad{}
\faGlobe\enspace{}\href{https://rusen.ai}{rusen.ai}\par\vspace{4pt}
\faGithub\enspace{}\href{https://github.com/rusenbb}{github.com/rusenbb}\quad{}
\faLinkedin\enspace{}\href{https://linkedin.com/in/rusenbirben}{linkedin.com/in/rusenbirben}
}\par\vspace{6pt}
\cvbody{İTÜ Yapay Zeka ve Veri Mühendisliği mezunu ve ODTÜ Uygulamalı NLP Grubu’nda bağımsız araştırmacıyım. Araştırma ilgim LLM’lerin gürbüzlüğü ve güvenilirliği; deneyimim Türkçe NLP, makine unutma ve finans alanında LLM uygulamalarını kapsıyor.}

\cvsection{\faGraduationCap}{EĞİTİM}
\begin{minipage}{\linewidth}\raggedright
\textbf{Yapay Zeka ve Veri Mühendisliği Lisansı}\dates{2022 - 2025}\par
{\small{}İstanbul Teknik Üniversitesi}\dates{ORT: 3.61 / 4.00}\par
\end{minipage}\par\vspace{5pt}
\begin{minipage}{\linewidth}\raggedright
\textbf{Elektronik ve Haberleşme Mühendisliği (lisans öğrenimi)}\dates{2020 - 2022}\par
{\small{}İstanbul Teknik Üniversitesi}\par
{\small{}\itshape{}\color{muted}Yapay Zeka ve Veri Mühendisliği'ne yatay geçiş yaptım.}\par\end{minipage}\par\vspace{5pt}


\cvsection{\faTrophy}{BURSLAR VE ÖDÜLLER}
\begin{minipage}{\linewidth}\raggedright
\textbf{\small{}1416 YLSY Bursu · Yüksek Lisans ve Doktora}\dates{2026}\par\vspace{2pt}
\cvbody{T.C. Millî Eğitim Bakanlığının yurt dışında yüksek lisans ve doktora eğitimi için verdiği 1416 YLSY bursunu kazandım.}
\end{minipage}\par\vspace{3pt}
\cvsection{\faFlask}{ARAŞTIRMA DENEYİMİ}
\begin{minipage}{\linewidth}\raggedright
\textbf{Bağımsız Araştırmacı}\dates{03/2026 - Bugün}\par
{\small{}\itshape{}\href{https://nlp.ceng.metu.edu.tr}{ODTÜ Uygulamalı NLP Grubu\extarrow}}\dates{Ankara, TR}\par\vspace{3pt}
\cvbody{Çağrı Toraman ile SIGTURK 2027 kapsamında Türk dilleri ve kültürel bağlamı değerlendiren benchmark için veri oluşturma ve değerlendirme metodolojisine katkı sağlıyor; LlamaTurk2 için eğitim verisi toplama ve ön eğitim deneylerinde çalışıyorum. Ayrıca \href{https://www.youtube.com/@metunlpconnect}{METU NLP Connect} kanalını yönetiyor ve uzman röportajları hazırlıyorum.}
\end{minipage}\par\vspace{6pt}
\begin{minipage}{\linewidth}\raggedright
\textbf{AIOps Araştırma Stajyeri}\dates{08/2024 - 09/2024}\par
{\small{}\itshape{}\href{https://havelsan.com/tr}{Havelsan\extarrow}}\dates{İstanbul, TR}\par\vspace{3pt}
\cvbody{Kestirimci bakım için yazılım loglarının analizini araştırdım; HDFS loglarıyla prototipler geliştirdim.}
\end{minipage}\par\vspace{6pt}
\begin{minipage}{\linewidth}\raggedright
\textbf{NLP Araştırma Stajyeri}\dates{07/2024 - 08/2024}\par
{\small{}\itshape{}\href{https://nlp.itu.edu.tr/}{İTÜ NLP Laboratuvarı\extarrow}}\dates{İstanbul, TR}\par\vspace{3pt}
\cvbody{Finans alanında LLM uygulamalarını araştırdım: duygu analizi, araç kullanımı, akıl yürütme, tahminleme ve bilgi getirme destekli üretim (RAG). Bu süreçte bitirme projem Kaira’ya katkı sağlayan deneylerin bir kısmını yürüttüm.}
\end{minipage}\par\vspace{6pt}
\cvsection{\faBriefcase}{PROFESYONEL DENEYİM}
\begin{minipage}{\linewidth}\raggedright
\textbf{Kurucu Ortak ve Teknik Lider}\dates{08/2025 - 04/2026}\par
{\small{}\itshape{}\href{https://github.com/fiction-ai-studios}{Fiction Studios\extarrow}}\dates{Ankara, TR}\par\vspace{3pt}
\cvbody{Bilkent Cyberpark’ta Fiction Studios’u bir arkadaşımla kurdum ve teknik liderliği üstlendim. Teknik geliştirmeyi yönettim; iki kurucuya destek veren üç arkadaşımızın çalışmalarını koordine ettim. Müzeler ve turistler için çok dilli, LLM destekli bir rehber uygulaması geliştirdik; sohbet tabanlı ziyaretçi deneyimi üzerine çalışırken müşteri ve yatırımcı görüşmeleri de yürüttük.}
\end{minipage}\par\vspace{6pt}
\begin{minipage}{\linewidth}\raggedright
\textbf{Junior Yapay Zeka Mühendisi}\dates{03/2025 - 07/2025}\par
{\small{}\itshape{}\href{https://cyberquote.com/tk/}{CyberQuote\extarrow}}\dates{İstanbul, TR}\par\vspace{3pt}
\cvbody{Phillip Capital iştiraki CyberQuote’un yeni kurulan yapay zeka ekibine ikinci üye olarak katıldım. Haberlerden duygu analizi çıkaran veri işleme akışları, finans uzmanları için kurum içi LLM araçları, web entegrasyonları ve mobil sesli asistan üzerinde çalıştım. Güncel finansal bağlamın dağınık kaynaklar ve PDF raporlarından edinilmesi, eksik bilginin güvenilmez yanıtlara yol açması gibi üretim sorunlarıyla uğraştım.}
\end{minipage}\par\vspace{6pt}

\newpage
\cvsection{\faCode}{SEÇİLMİŞ PROJELER}
\begin{minipage}{\linewidth}\raggedright
{\ttfamily{}\bfseries{}Kaira}\dates{2025}\par
{\small{}\itshape{}\color{muted}Bitirme Projesi: Finansal Akıl Yürütme Modeli}\hfill{\small{}%
%
}\par\vspace{4pt}
\cvbody{İTÜ bitirme projemde Qwen tabanlı dil modellerinin finansal akıl yürütme becerilerini araştırdım. Unsloth ve LoRA kullanarak GRPO ile ince ayar yaptım; FinQA ve GSM8K üzerinde ödül tasarımı ve akıl yürütme davranışını inceleyen deneyler yürüttüm. Ollama tabanlı bir sohbet uygulaması da geliştirdim.}
\end{minipage}\par\vspace{5pt}
\begin{minipage}{\linewidth}\raggedright
{\ttfamily{}\bfseries{}AIZheimer}\dates{2024}\par
{\small{}\itshape{}\color{muted}Stable Diffusion’da Makine Unutma}\hfill{\small{}%
\href{https://github.com/rusenbb/ai-zheimer}{\faGithub\enspace{}Kod}%
}\par\vspace{4pt}
\cvbody{Forget-Me-Not çalışmasını temel alarak Stable Diffusion 2.1 için çapraz dikkat tabanlı makine unutma yöntemlerini uyguladım ve uyarladım. Attention controller/processor kodu geliştirerek attention re-steering, distraction ve random stimulus kayıplarını inceledim. Hedef kavramı unutturma ile diğer üretim yeteneklerini koruma arasındaki dengeyi nitel deneyler ve hazırladığım görselleştirmelerle değerlendirdim.}
\end{minipage}\par\vspace{5pt}
\begin{minipage}{\linewidth}\raggedright
{\ttfamily{}\bfseries{}Biting The Bytes}\dates{2024}\par
{\small{}\itshape{}\color{muted}Türkçe Karakter Tamamlama}\hfill{\small{}%
\href{https://github.com/rusen-archive/Biting-The-Bytes}{\faGithub\enspace{}Kod}\quad{}\href{https://huggingface.co/spaces/rusen/diacritize-tr}{\faExternalLink\enspace{}Demo}%
}\par\vspace{4pt}
\cvbody{Türkçe karakter restorasyonunu bayt düzeyinde token sınıflandırması olarak ele alarak ByT5-small modelini PEFT/LoRA ile eğittim. Model eğitimi ve hiperparametre ayarlamasını ağırlıklı olarak üstlendim; hata ayıklama ve Hugging Face Spaces’te yayımlama sürecini ekip arkadaşımla paylaştım. Proje, gürültülü veya yalnızca ASCII karakterleri içeren Türkçe metinlerde ş, ç, ı, ğ, ü ve ö karakterlerinin geri kazandırılmasına odaklandı.}
\end{minipage}\par\vspace{5pt}
\begin{minipage}{\linewidth}\raggedright
{\ttfamily{}\bfseries{}To-AI-or-Not-to-AI}\dates{2023}\par
{\small{}\itshape{}\color{muted}Yapay Zeka Metin Tespiti}\hfill{\small{}%
\href{https://github.com/rusen-archive/To-AI-or-Not-to-AI}{\faGithub\enspace{}Kod}\quad{}\href{https://huggingface.co/spaces/rusen/gpt_detector}{\faExternalLink\enspace{}Demo}%
}\par\vspace{4pt}
\cvbody{Üç RoBERTa tabanlı dedektörü ince ayarladım ve tahminlerini çoğunluk oylamasıyla birleştirdim. 26.270 örnekten oluşan veri kümesinin oluşturulmasını ve ön işlemesini ekip arkadaşımla paylaştım; uygulamayı Hugging Face Spaces’te yayımladım. Topluluk modeli, aynı dağılımdan rastgele ayrılmış test kümesinde \%94,21 doğruluk sağladı; farklı diller ve belirgin biçimde farklı alanlar bu değerlendirmeye dahil değildi.}
\end{minipage}\par\vspace{5pt}
\begin{minipage}{\linewidth}\raggedright
{\ttfamily{}\bfseries{}Tobor}\dates{2024 - 2025}\par
{\small{}\itshape{}\color{muted}Otonom Depo Robotu}\hfill{\small{}%
\href{https://github.com/rusen-archive/Amazon-Warehouse-Robot}{\faGithub\enspace{}Kod}\quad{}\href{https://youtu.be/h-Jxs_y9kNA}{\faYoutubePlay\enspace{}Video}%
}\par\vspace{4pt}
\cvbody{ROS2 üzerinde navigasyon ve YOLO tabanlı algıyı birleştiren otonom depo robotunu ekip olarak geliştirdik. Kutuları alma ve bırakma mekanizmasını yaptım; veri etiketlemeyi ekiple birlikte yürüttüm ve diğer algoritmaların geliştirilmesine destek verdim.}
\end{minipage}\par\vspace{5pt}

\cvsection{\faCogs}{TEKNİK BECERİLER}
{\small{}\textbf{Diller}\enspace{} Python / JavaScript/TypeScript / SQL / C/C++}\par\vspace{4pt}
{\small{}\textbf{Çerçeveler ve Araçlar}\enspace{} PyTorch / TensorFlow / Transformers / PEFT / LoRA / Unsloth / Scikit-learn / Pandas / NumPy / ROS2 / YOLO}\par\vspace{4pt}
{\small{}\textbf{Uzmanlık Alanları}\enspace{} NLP / LLM’ler / RAG / Makine Unutma / Model Değerlendirme}\par\vspace{4pt}

\cvsection{\faCertificate}{KURSLAR / DİLLER}
\textbf{\small{}\href{https://coursera.org/share/9c59e9758d378fb8a8b4ab7074dd71ed}{Deep Learning Specialization\extarrow}}\dates{DeepLearning.AI · Coursera}\par\vspace{2pt}
\cvbody{Sinir ağları, CNN’ler, sıralı modeller, optimizasyon ve makine öğrenmesi projelerinin yapılandırılması.}\vspace{5pt}
\textbf{\small{}\href{https://coursera.org/share/f8ab2981ab0bbdc3fec17c6e88804ff8}{AWS Cloud Technical Essentials\extarrow}}\dates{AWS · Coursera}\par\vspace{2pt}
\cvbody{Bulut mimarisi; EC2, Lambda, S3, RDS ve DynamoDB dahil AWS servisleri.}\vspace{5pt}

\vspace{5pt}
\faLanguage\enspace{}
{\small{}\textbf{Türkçe}\enspace{}Anadil\quad{}\textbf{İngilizce}\enspace{}TOEFL iBT 109/120 (8 Kasım 2025)}\par

\cvsection{\faHeartO}{İLGİ ALANLARI\hfill{\normalfont{}\small{}\href{https://rusenbirben.com}{rusenbirben.com\extarrow}}}
\noindent\begin{minipage}[t]{0.235\linewidth}\raggedright
{\large{}\faCoffee}\par\vspace{5pt}{\ttfamily{}\bfseries{}\small{}Kahve}\par\vspace{4pt}
\cvbody{V60 / Aeropress / Moka Pot}
\end{minipage}%
\hfill%
\noindent\begin{minipage}[t]{0.235\linewidth}\raggedright
{\large{}\faBook}\par\vspace{5pt}{\ttfamily{}\bfseries{}\small{}Kitap}\par\vspace{4pt}
\cvbody{Felsefe / YZ/Bilgisayar / Finans}
\end{minipage}%
\hfill%
\noindent\begin{minipage}[t]{0.235\linewidth}\raggedright
{\large{}\faCamera}\par\vspace{5pt}{\ttfamily{}\bfseries{}\small{}Fotoğrafçılık}\par\vspace{4pt}
\cvbody{Sokak fotoğrafçılığı / Çift pozlama / Kuş fotoğrafçılığı}
\end{minipage}%
\hfill%
\noindent\begin{minipage}[t]{0.235\linewidth}\raggedright
{\large{}\faFilm}\par\vspace{5pt}{\ttfamily{}\bfseries{}\small{}Anime ve Manga}\par\vspace{4pt}
\cvbody{İzleyici / Okuyucu / Koleksiyoncu}
\end{minipage}%
\par
\end{document}
